\documentclass[11pt]{article}
\usepackage{amsmath, amssymb, amsthm, geometry, graphicx, enumitem, xcolor}
\geometry{a4paper, margin=1in}

\title{\textbf{Homework 4.6: Change of Basis}}
\author{Tashfeen Omran}
\date{September 2026}

\begin{document}

\maketitle
\tableofcontents
\newpage

% ==========================================
% PART 1: INTRODUCTION
% ==========================================
\section{Part 1: Comprehensive Introduction, Context, and Prerequisites}

\subsection{Core Concepts}
In linear algebra, a \textbf{basis} for a vector space $V$ is a set of linearly independent vectors that spans $V$. You can think of a basis as a ``coordinate system'' or a ``language'' used to describe vectors. If we have a vector $\mathbf{x}$ in $V$ and a basis $\mathcal{B} = \{\mathbf{b}_1, \mathbf{b}_2, \dots, \mathbf{b}_n\}$, we can uniquely express $\mathbf{x}$ as a linear combination of the basis vectors:
\[ \mathbf{x} = c_1\mathbf{b}_1 + c_2\mathbf{b}_2 + \dots + c_n\mathbf{b}_n \]
The weights $c_1, \dots, c_n$ are called the \textbf{coordinates of $\mathbf{x}$ relative to $\mathcal{B}$}, and they form the coordinate vector $[\mathbf{x}]_{\mathcal{B}}$.

The \textbf{Change of Basis} problem asks: If we know the coordinates of a vector in one basis $\mathcal{B}$, how do we find its coordinates in a different basis $\mathcal{C}$? Because the mapping from $\mathcal{B}$-coordinates to $\mathcal{C}$-coordinates is a linear transformation, it can be represented by matrix multiplication. This matrix is called the \textbf{change-of-coordinates matrix}, denoted as $P_{\mathcal{C} \leftarrow \mathcal{B}}$.

\subsection{Intuition and Derivation}
Imagine you are measuring a physical object. Basis $\mathcal{B}$ might represent measurements in inches, while basis $\mathcal{C}$ represents measurements in centimeters. The object itself (the vector $\mathbf{x}$) does not change, but its numerical representation (the coordinate vector) changes depending on the system you use. 

To derive the matrix $P_{\mathcal{C} \leftarrow \mathcal{B}}$, we apply the coordinate mapping to the basis vectors of $\mathcal{B}$. Since $\mathbf{x} = c_1\mathbf{b}_1 + \dots + c_n\mathbf{b}_n$, taking the $\mathcal{C}$-coordinates of both sides yields:
\[ [\mathbf{x}]_{\mathcal{C}} = c_1[\mathbf{b}_1]_{\mathcal{C}} + \dots + c_n[\mathbf{b}_n]_{\mathcal{C}} \]
By the definition of matrix-vector multiplication, this is equivalent to:
\[ [\mathbf{x}]_{\mathcal{C}} = \begin{bmatrix} [\mathbf{b}_1]_{\mathcal{C}} & [\mathbf{b}_2]_{\mathcal{C}} & \dots & [\mathbf{b}_n]_{\mathcal{C}} \end{bmatrix} \begin{bmatrix} c_1 \\ c_2 \\ \vdots \\ c_n \end{bmatrix} \]
Thus, the change-of-coordinates matrix $P_{\mathcal{C} \leftarrow \mathcal{B}}$ is simply the matrix whose columns are the $\mathcal{C}$-coordinate vectors of the basis vectors in $\mathcal{B}$.

\subsection{Historical Context and Motivation}
The concept of coordinates was introduced by René Descartes in the 17th century, bridging geometry and algebra. However, the formalization of abstract vector spaces and the realization that multiple coordinate systems (bases) could describe the same space emerged in the late 19th and early 20th centuries through the work of mathematicians like Giuseppe Peano and Hermann Grassmann. The motivation was to simplify complex problems: a problem that is computationally difficult in the standard basis might become trivial if transformed into a specialized basis (such as an eigenbasis, which diagonalizes a matrix).

\subsection{Key Formulas}
\begin{itemize}
    \item \textbf{The Change of Coordinates Equation:} 
    \[ [\mathbf{x}]_{\mathcal{C}} = P_{\mathcal{C} \leftarrow \mathcal{B}} [\mathbf{x}]_{\mathcal{B}} \]
    \item \textbf{Construction of the Matrix:} 
    \[ P_{\mathcal{C} \leftarrow \mathcal{B}} = \begin{bmatrix} [\mathbf{b}_1]_{\mathcal{C}} & [\mathbf{b}_2]_{\mathcal{C}} & \dots & [\mathbf{b}_n]_{\mathcal{C}} \end{bmatrix} \]
    \item \textbf{Invertibility:} 
    \[ (P_{\mathcal{C} \leftarrow \mathcal{B}})^{-1} = P_{\mathcal{B} \leftarrow \mathcal{C}} \]
    \item \textbf{Row Reduction Shortcut (for $\mathbb{R}^n$):} 
    \[ \begin{bmatrix} \mathbf{c}_1 & \dots & \mathbf{c}_n \mid \mathbf{b}_1 & \dots & \mathbf{b}_n \end{bmatrix} \sim \begin{bmatrix} I \mid P_{\mathcal{C} \leftarrow \mathcal{B}} \end{bmatrix} \]
\end{itemize}

\subsection{Prerequisites}
To master this topic, you must be proficient in:
1. \textbf{Row Reduction (Gaussian Elimination):} Transforming augmented matrices into Reduced Row Echelon Form (RREF).
2. \textbf{Matrix Multiplication:} Understanding how a matrix acts on a vector as a linear combination of its columns.
3. \textbf{Vector Space Axioms:} Understanding span, linear independence, and the definition of a basis.
4. \textbf{Isomorphisms:} Recognizing that abstract spaces (like polynomials $\mathbb{P}_n$) behave exactly like $\mathbb{R}^{n+1}$ when coordinates are used.

% ==========================================
% PART 2: DETAILED HOMEWORK SOLUTIONS
% ==========================================
\section{Part 2: Detailed Homework Solutions}
\textit{Note: To maximize learning efficiency, a strategic selection of 8 problems has been chosen, covering at least two of every distinct problem archetype found in the text.}

\subsection{Type 1: Nonstandard Bases in $\mathbb{R}^n$}

\textbf{Problem 1:} Let $\mathcal{B} = \{\mathbf{b}_1, \mathbf{b}_2\}$ and $\mathcal{C} = \{\mathbf{c}_1, \mathbf{c}_2\}$ be bases for a vector space $V$, and suppose $\mathbf{b}_1 = 6\mathbf{c}_1 - 2\mathbf{c}_2$ and $\mathbf{b}_2 = 9\mathbf{c}_1 - 4\mathbf{c}_2$.
(a) Find the change-of-coordinates matrix from $\mathcal{B}$ to $\mathcal{C}$.
(b) Find $[\mathbf{x}]_{\mathcal{C}}$ for $\mathbf{x} = -3\mathbf{b}_1 + 2\mathbf{b}_2$.

\textbf{Solution:}
\textbf{Part (a):} The change-of-coordinates matrix $P_{\mathcal{C} \leftarrow \mathcal{B}}$ is formed by placing the $\mathcal{C}$-coordinate vectors of the $\mathcal{B}$ basis vectors into columns.
From the given equations, we can extract the coordinates directly:
\[ [\mathbf{b}_1]_{\mathcal{C}} = \begin{bmatrix} 6 \\ -2 \end{bmatrix}, \quad [\mathbf{b}_2]_{\mathcal{C}} = \begin{bmatrix} 9 \\ -4 \end{bmatrix} \]
Therefore, the matrix is:
\[ P_{\mathcal{C} \leftarrow \mathcal{B}} = \begin{bmatrix} 6 & 9 \\ -2 & -4 \end{bmatrix} \]

\textbf{Part (b):} We are given $\mathbf{x} = -3\mathbf{b}_1 + 2\mathbf{b}_2$. This means the $\mathcal{B}$-coordinate vector of $\mathbf{x}$ is:
\[ [\mathbf{x}]_{\mathcal{B}} = \begin{bmatrix} -3 \\ 2 \end{bmatrix} \]
To find $[\mathbf{x}]_{\mathcal{C}}$, we use the change of coordinates equation:
\[ [\mathbf{x}]_{\mathcal{C}} = P_{\mathcal{C} \leftarrow \mathcal{B}} [\mathbf{x}]_{\mathcal{B}} \]
\[ [\mathbf{x}]_{\mathcal{C}} = \begin{bmatrix} 6 & 9 \\ -2 & -4 \end{bmatrix} \begin{bmatrix} -3 \\ 2 \end{bmatrix} = \begin{bmatrix} (6)(-3) + (9)(2) \\ (-2)(-3) + (-4)(2) \end{bmatrix} = \begin{bmatrix} -18 + 18 \\ 6 - 8 \end{bmatrix} = \begin{bmatrix} 0 \\ -2 \end{bmatrix} \]
\textbf{Final Answer:} (a) $\begin{bmatrix} 6 & 9 \\ -2 & -4 \end{bmatrix}$, (b) $\begin{bmatrix} 0 \\ -2 \end{bmatrix}$.

\vspace{0.5cm}

\textbf{Problem 7:} Let $\mathbf{b}_1 = \begin{bmatrix} 7 \\ 5 \end{bmatrix}, \mathbf{b}_2 = \begin{bmatrix} -3 \\ -1 \end{bmatrix}, \mathbf{c}_1 = \begin{bmatrix} 1 \\ -5 \end{bmatrix}, \mathbf{c}_2 = \begin{bmatrix} -2 \\ 2 \end{bmatrix}$. Find the change-of-coordinates matrix from $\mathcal{B}$ to $\mathcal{C}$.

\textbf{Solution:}
We need to find $P_{\mathcal{C} \leftarrow \mathcal{B}}$. We use the augmented matrix method: $[ \mathcal{C} \mid \mathcal{B} ] \sim [ I \mid P_{\mathcal{C} \leftarrow \mathcal{B}} ]$.
Set up the augmented matrix:
\[ \begin{bmatrix} \mathbf{c}_1 & \mathbf{c}_2 \mid \mathbf{b}_1 & \mathbf{b}_2 \end{bmatrix} = \begin{bmatrix} 1 & -2 & 7 & -3 \\ -5 & 2 & 5 & -1 \end{bmatrix} \]
Row reduce to RREF. First, eliminate the $-5$ in the second row ($R_2 \leftarrow R_2 + 5R_1$):
\[ \begin{bmatrix} 1 & -2 & 7 & -3 \\ 0 & -8 & 40 & -16 \end{bmatrix} \]
Divide the second row by $-8$ ($R_2 \leftarrow -\frac{1}{8}R_2$):
\[ \begin{bmatrix} 1 & -2 & 7 & -3 \\ 0 & 1 & -5 & 2 \end{bmatrix} \]
Eliminate the $-2$ in the first row ($R_1 \leftarrow R_1 + 2R_2$):
\[ \begin{bmatrix} 1 & 0 & -3 & 1 \\ 0 & 1 & -5 & 2 \end{bmatrix} \]
The right half of the matrix is $P_{\mathcal{C} \leftarrow \mathcal{B}}$.
\textbf{Final Answer:} $P_{\mathcal{C} \leftarrow \mathcal{B}} = \begin{bmatrix} -3 & 1 \\ -5 & 2 \end{bmatrix}$.

\subsection{Type 2: Abstract Vector Spaces}

\textbf{Problem 15:} In $\mathbb{P}_2$, find the change-of-coordinates matrix from the basis $\mathcal{B} = \{1-2t+t^2, 3-5t+4t^2, 2t+3t^2\}$ to the standard basis $\mathcal{C} = \{1, t, t^2\}$. Then find the $\mathcal{B}$-coordinate vector for $-1+2t$.

\textbf{Solution:}
First, find $P_{\mathcal{C} \leftarrow \mathcal{B}}$. Since $\mathcal{C}$ is the standard basis, the $\mathcal{C}$-coordinates of a polynomial $a + bt + ct^2$ are simply $\begin{bmatrix} a \\ b \\ c \end{bmatrix}$.
\[ [\mathbf{b}_1]_{\mathcal{C}} = \begin{bmatrix} 1 \\ -2 \\ 1 \end{bmatrix}, \quad [\mathbf{b}_2]_{\mathcal{C}} = \begin{bmatrix} 3 \\ -5 \\ 4 \end{bmatrix}, \quad [\mathbf{b}_3]_{\mathcal{C}} = \begin{bmatrix} 0 \\ 2 \\ 3 \end{bmatrix} \]
Thus, the change-of-coordinates matrix is:
\[ P_{\mathcal{C} \leftarrow \mathcal{B}} = \begin{bmatrix} 1 & 3 & 0 \\ -2 & -5 & 2 \\ 1 & 4 & 3 \end{bmatrix} \]
Next, we need to find $[\mathbf{x}]_{\mathcal{B}}$ for $\mathbf{x} = -1 + 2t$. The $\mathcal{C}$-coordinate vector for $\mathbf{x}$ is:
\[ [\mathbf{x}]_{\mathcal{C}} = \begin{bmatrix} -1 \\ 2 \\ 0 \end{bmatrix} \]
We must solve the equation $P_{\mathcal{C} \leftarrow \mathcal{B}} [\mathbf{x}]_{\mathcal{B}} = [\mathbf{x}]_{\mathcal{C}}$. We set up the augmented matrix:
\[ \begin{bmatrix} 1 & 3 & 0 & -1 \\ -2 & -5 & 2 & 2 \\ 1 & 4 & 3 & 0 \end{bmatrix} \]
Row reduce to RREF:
$R_2 \leftarrow R_2 + 2R_1$:
\[ \begin{bmatrix} 1 & 3 & 0 & -1 \\ 0 & 1 & 2 & 0 \\ 1 & 4 & 3 & 0 \end{bmatrix} \]
$R_3 \leftarrow R_3 - R_1$:
\[ \begin{bmatrix} 1 & 3 & 0 & -1 \\ 0 & 1 & 2 & 0 \\ 0 & 1 & 3 & 1 \end{bmatrix} \]
$R_3 \leftarrow R_3 - R_2$:
\[ \begin{bmatrix} 1 & 3 & 0 & -1 \\ 0 & 1 & 2 & 0 \\ 0 & 0 & 1 & 1 \end{bmatrix} \]
$R_2 \leftarrow R_2 - 2R_3$:
\[ \begin{bmatrix} 1 & 3 & 0 & -1 \\ 0 & 1 & 0 & -2 \\ 0 & 0 & 1 & 1 \end{bmatrix} \]
$R_1 \leftarrow R_1 - 3R_2$:
\[ \begin{bmatrix} 1 & 0 & 0 & 5 \\ 0 & 1 & 0 & -2 \\ 0 & 0 & 1 & 1 \end{bmatrix} \]
The rightmost column is $[\mathbf{x}]_{\mathcal{B}}$.
\textbf{Final Answer:} $P_{\mathcal{C} \leftarrow \mathcal{B}} = \begin{bmatrix} 1 & 3 & 0 \\ -2 & -5 & 2 \\ 1 & 4 & 3 \end{bmatrix}$ and $[\mathbf{x}]_{\mathcal{B}} = \begin{bmatrix} 5 \\ -2 \\ 1 \end{bmatrix}$.

\vspace{0.5cm}

\textbf{Problem 19:} Let $\mathcal{B} = \{x_0, \dots, x_6\}$ and $\mathcal{C} = \{y_0, \dots, y_6\}$, where $x_k$ is the function $\cos^k t$ and $y_k$ is the function $\cos kt$. (a) Set $P = \begin{bmatrix} [y_0]_{\mathcal{B}} & \dots & [y_6]_{\mathcal{B}} \end{bmatrix}$ and calculate $P$. (b) Explain why the columns of $P^{-1}$ are the $\mathcal{C}$-coordinate vectors of $x_0, \dots, x_6$.

\textbf{Solution:}
\textbf{Part (a):} We use the trigonometric identities provided in the text (Exercise 53/54) to express each $y_k$ in terms of the basis $\mathcal{B}$ (powers of cosine).
\begin{align*}
y_0 &= \cos 0t = 1 = 1x_0 \\
y_1 &= \cos 1t = \cos t = 1x_1 \\
y_2 &= \cos 2t = -1 + 2\cos^2 t = -1x_0 + 2x_2 \\
y_3 &= \cos 3t = -3\cos t + 4\cos^3 t = -3x_1 + 4x_3 \\
y_4 &= \cos 4t = 1 - 8\cos^2 t + 8\cos^4 t = 1x_0 - 8x_2 + 8x_4 \\
y_5 &= \cos 5t = 5\cos t - 20\cos^3 t + 16\cos^5 t = 5x_1 - 20x_3 + 16x_5 \\
y_6 &= \cos 6t = -1 + 18\cos^2 t - 48\cos^4 t + 32\cos^6 t = -1x_0 + 18x_2 - 48x_4 + 32x_6
\end{align*}
Extracting the coefficients to form the coordinate vectors $[y_k]_{\mathcal{B}}$ yields the columns of $P$:
\[ P = \begin{bmatrix} 
1 & 0 & -1 & 0 & 1 & 0 & -1 \\
0 & 1 & 0 & -3 & 0 & 5 & 0 \\
0 & 0 & 2 & 0 & -8 & 0 & 18 \\
0 & 0 & 0 & 4 & 0 & -20 & 0 \\
0 & 0 & 0 & 0 & 8 & 0 & -48 \\
0 & 0 & 0 & 0 & 0 & 16 & 0 \\
0 & 0 & 0 & 0 & 0 & 0 & 32 
\end{bmatrix} \]

\textbf{Part (b):} By definition, the matrix $P$ constructed in part (a) is the change-of-coordinates matrix from $\mathcal{C}$ to $\mathcal{B}$, denoted $P_{\mathcal{B} \leftarrow \mathcal{C}}$. A fundamental theorem of change of basis states that the inverse of $P_{\mathcal{B} \leftarrow \mathcal{C}}$ is $P_{\mathcal{C} \leftarrow \mathcal{B}}$. Therefore, $P^{-1} = P_{\mathcal{C} \leftarrow \mathcal{B}}$. By definition, the columns of $P_{\mathcal{C} \leftarrow \mathcal{B}}$ are the $\mathcal{C}$-coordinate vectors of the basis vectors in $\mathcal{B}$ (which are $x_0, \dots, x_6$).
\textbf{Final Answer:} Matrix $P$ is shown above. $P^{-1} = P_{\mathcal{C} \leftarrow \mathcal{B}}$, whose columns are $[x_k]_{\mathcal{C}}$.

\subsection{Type 3: Conceptual True/False}

\textbf{Problem 11 \& 12:} Mark each statement True or False. Justify.
11. The columns of the change-of-coordinates matrix $P_{\mathcal{C} \leftarrow \mathcal{B}}$ are $\mathcal{B}$-coordinate vectors of the vectors in $\mathcal{C}$.
12. The columns of $P_{\mathcal{C} \leftarrow \mathcal{B}}$ are linearly independent.

\textbf{Solution:}
\textbf{11. False.} The columns of $P_{\mathcal{C} \leftarrow \mathcal{B}}$ are the $\mathcal{C}$-coordinate vectors of the vectors in $\mathcal{B}$, not the other way around. The notation $P_{\mathcal{C} \leftarrow \mathcal{B}}$ means we are converting \textit{from} $\mathcal{B}$ \textit{to} $\mathcal{C}$, so we must evaluate the $\mathcal{B}$ vectors in the language of $\mathcal{C}$.
\textbf{12. True.} The columns of $P_{\mathcal{C} \leftarrow \mathcal{B}}$ are the coordinate vectors of a basis ($\mathcal{B}$). Because the vectors in $\mathcal{B}$ are linearly independent, their coordinate vectors in any basis $\mathcal{C}$ are also linearly independent. Furthermore, $P_{\mathcal{C} \leftarrow \mathcal{B}}$ is an invertible matrix, and all invertible matrices have linearly independent columns.

\subsection{Type 4: Proofs and Derivations}

\textbf{Problem 17:} Given $\mathbf{v}$ in $V$, there exist scalars $x_1, \dots, x_n$ such that $\mathbf{v} = x_1\mathbf{b}_1 + x_2\mathbf{b}_2 + \dots + x_n\mathbf{b}_n$ because (a) \_\_\_\_\_\_\_\_\_\_\_\_. Apply the coordinate mapping determined by the basis $\mathcal{C}$, and obtain $[\mathbf{v}]_{\mathcal{C}} = x_1[\mathbf{b}_1]_{\mathcal{C}} + \dots + x_n[\mathbf{b}_n]_{\mathcal{C}}$ because (b) \_\_\_\_\_\_\_\_\_\_\_\_. This equation may be written in the form $[\mathbf{v}]_{\mathcal{C}} = \begin{bmatrix} [\mathbf{b}_1]_{\mathcal{C}} & \dots & [\mathbf{b}_n]_{\mathcal{C}} \end{bmatrix} \begin{bmatrix} x_1 \\ \vdots \\ x_n \end{bmatrix}$ by the definition of (c) \_\_\_\_\_\_\_\_\_\_\_\_. This shows that the matrix $P_{\mathcal{C} \leftarrow \mathcal{B}}$ satisfies $[\mathbf{v}]_{\mathcal{C}} = P_{\mathcal{C} \leftarrow \mathcal{B}}[\mathbf{v}]_{\mathcal{B}}$ for each $\mathbf{v}$ in $V$, because the vector on the right side of the equation is (d) \_\_\_\_\_\_\_\_\_\_\_\_.

\textbf{Solution:}
(a) $\mathcal{B}$ is a basis for $V$, meaning it spans $V$, so any vector $\mathbf{v}$ can be written as a linear combination of the basis vectors.
(b) The coordinate mapping $\mathbf{x} \mapsto [\mathbf{x}]_{\mathcal{C}}$ is a linear transformation, which preserves vector addition and scalar multiplication.
(c) Matrix-vector multiplication (a matrix times a vector is a linear combination of the columns of the matrix).
(d) $[\mathbf{v}]_{\mathcal{B}}$, which is the coordinate vector of $\mathbf{v}$ relative to $\mathcal{B}$.

\textbf{Problem 18:} Suppose $Q$ is any matrix such that $[\mathbf{v}]_{\mathcal{C}} = Q[\mathbf{v}]_{\mathcal{B}}$ for each $\mathbf{v}$ in $V$. Set $\mathbf{v} = \mathbf{b}_1$. Then the equation shows that $[\mathbf{b}_1]_{\mathcal{C}}$ is the first column of $Q$ because (a) \_\_\_\_\_\_\_\_\_\_\_\_. Similarly, for $k = 2, \dots, n$, the $k$-th column of $Q$ is (b) \_\_\_\_\_\_\_\_\_\_\_\_ because (c) \_\_\_\_\_\_\_\_\_\_\_\_. This shows that the matrix $P_{\mathcal{C} \leftarrow \mathcal{B}}$ is the unique matrix that satisfies the condition.

\textbf{Solution:}
(a) $[\mathbf{b}_1]_{\mathcal{B}} = \mathbf{e}_1 = \begin{bmatrix} 1 \\ 0 \\ \vdots \\ 0 \end{bmatrix}$. Multiplying $Q$ by $\mathbf{e}_1$ extracts exactly the first column of $Q$.
(b) $[\mathbf{b}_k]_{\mathcal{C}}$.
(c) $[\mathbf{b}_k]_{\mathcal{B}} = \mathbf{e}_k$, and multiplying $Q$ by $\mathbf{e}_k$ extracts the $k$-th column of $Q$.

% ==========================================
% PART 3: IN-DEPTH ANALYSIS
% ==========================================
\section{Part 3: In-Depth Analysis of Problems and Techniques}

\subsection{A) Problem Types and General Approach}
\begin{itemize}
    \item \textbf{Direct Matrix Construction (e.g., Problem 1):} You are given the explicit linear combinations of how one basis is formed from another. \textit{Strategy:} Extract the coefficients directly to form the columns of the change-of-basis matrix.
    \item \textbf{Row Reduction Method (e.g., Problem 7):} You are given two nonstandard bases in $\mathbb{R}^n$. \textit{Strategy:} Set up the augmented matrix $[ \mathcal{C} \mid \mathcal{B} ]$ and row reduce to $[ I \mid P_{\mathcal{C} \leftarrow \mathcal{B}} ]$.
    \item \textbf{Abstract Vector Spaces (e.g., Problems 15, 19):} The vectors are polynomials or functions, not column vectors. \textit{Strategy:} Map the abstract vectors to $\mathbb{R}^n$ using a standard basis (like $\{1, t, t^2\}$), then apply the standard matrix techniques.
    \item \textbf{Theoretical Proofs (e.g., Problems 17, 18):} Testing the foundational definitions. \textit{Strategy:} Rely heavily on the definition of linear transformations and matrix-vector multiplication.
\end{itemize}

\subsection{B) Key Algebraic and Calculus Manipulations}
\begin{itemize}
    \item \textbf{Augmented Matrix Row Reduction:} Used in Problem 7 and 15. This is the workhorse of linear algebra. It simultaneously solves $n$ systems of linear equations. By reducing $[ \mathcal{C} \mid \mathcal{B} ]$, we are solving $C\mathbf{x} = \mathbf{b}_i$ for every basis vector in $\mathcal{B}$.
    \item \textbf{Isomorphic Mapping:} In Problem 15, the polynomial $a + bt + ct^2$ was treated exactly like the column vector $\begin{bmatrix} a \\ b \\ c \end{bmatrix}$. This trick is crucial because it allows us to use matrix algebra on non-matrix objects.
    \item \textbf{Trigonometric Substitution:} In Problem 19, complex trigonometric identities were used to map functions of multiple angles ($\cos kt$) into powers of a single angle ($\cos^k t$). This allowed us to build a discrete matrix out of continuous functions.
\end{itemize}

% ==========================================
% PART 4: CHEATSHEET AND TIPS
% ==========================================
\section{Part 4: ``Cheatsheet'' and Tips for Success}

\subsection{Summary of Important Formulas}
\begin{itemize}
    \item $P_{\mathcal{C} \leftarrow \mathcal{B}} = \begin{bmatrix} [\mathbf{b}_1]_{\mathcal{C}} & \dots & [\mathbf{b}_n]_{\mathcal{C}} \end{bmatrix}$
    \item $[\mathbf{x}]_{\mathcal{C}} = P_{\mathcal{C} \leftarrow \mathcal{B}} [\mathbf{x}]_{\mathcal{B}}$
    \item $P_{\mathcal{B} \leftarrow \mathcal{C}} = (P_{\mathcal{C} \leftarrow \mathcal{B}})^{-1}$
\end{itemize}

\subsection{Cheats and Shortcuts}
\begin{itemize}
    \item \textbf{The $[ \mathcal{C} \mid \mathcal{B} ]$ Trick:} To find $P_{\mathcal{C} \leftarrow \mathcal{B}}$, put the ``target'' basis $\mathcal{C}$ on the left of the augmented matrix, and the ``source'' basis $\mathcal{B}$ on the right. Row reduce. The right side becomes your answer.
    \item \textbf{Standard Basis Shortcut:} If $\mathcal{C}$ is the standard basis (usually denoted $\mathcal{E}$), then $P_{\mathcal{E} \leftarrow \mathcal{B}}$ is simply the matrix whose columns are the vectors of $\mathcal{B}$ themselves. No row reduction needed!
\end{itemize}

\subsection{Common Pitfalls and Traps}
\begin{itemize}
    \item \textbf{Reversing the Direction:} The most common mistake is calculating $P_{\mathcal{B} \leftarrow \mathcal{C}}$ when asked for $P_{\mathcal{C} \leftarrow \mathcal{B}}$. Always remember: the basis you are converting \textit{to} goes on the left of the augmented matrix.
    \item \textbf{Multiplying the Wrong Vector:} If you have $[\mathbf{x}]_{\mathcal{C}}$ and you multiply it by $P_{\mathcal{C} \leftarrow \mathcal{B}}$, you will get garbage. The matrix $P_{\mathcal{C} \leftarrow \mathcal{B}}$ must consume a $\mathcal{B}$-vector to produce a $\mathcal{C}$-vector.
\end{itemize}

% ==========================================
% PART 5: CONCEPTUAL SYNTHESIS
% ==========================================
\section{Part 5: Conceptual Synthesis and The ``Big Picture''}

\subsection{A) Thematic Connections}
The core theme of this topic is \textbf{Invariance under Transformation}. A vector $\mathbf{x}$ is a geometric object that exists independently of how we measure it. The coordinate vector $[\mathbf{x}]_{\mathcal{B}}$ is just a numerical representation. Changing the basis is akin to changing units in physics (e.g., Celsius to Fahrenheit) or changing perspective in a 3D engine (World Space to Camera Space). The underlying reality remains unchanged, but the numbers shift according to a strict linear rule.

\subsection{B) Forward and Backward Links}
\begin{itemize}
    \item \textbf{Backward Links:} This topic relies heavily on the concepts of linear independence, span, and the Invertible Matrix Theorem. If a set of vectors doesn't form a basis, the change-of-coordinates matrix won't be invertible.
    \item \textbf{Forward Links:} This is the absolute bedrock for \textbf{Diagonalization} and \textbf{Eigenvalues}. Later, you will learn that for a matrix $A$, we can find a basis of eigenvectors $\mathcal{B}$. The matrix $P$ used to diagonalize $A$ ($A = PDP^{-1}$) is exactly the change-of-basis matrix $P_{\mathcal{E} \leftarrow \mathcal{B}}$.
\end{itemize}

% ==========================================
% PART 6: ABSTRACT REAL-WORLD MODELING
% ==========================================
\section{Part 6: Abstract Real-World Modeling (AI, ML, \& Quant Finance)}

\subsection{A) Concrete Scenarios}
1. \textbf{Principal Component Analysis (PCA) in Machine Learning:} In ML, high-dimensional data often contains redundant features. PCA finds a new orthogonal basis (the principal components) where the data's variance is maximized along the new axes. The transformation from the original feature space to the PCA space is a pure Change of Basis operation.
2. \textbf{Statistical Arbitrage in Quantitative Finance:} Quant funds trade baskets of correlated stocks. By changing the basis from individual stocks to ``Principal Portfolios'' (eigenvectors of the covariance matrix), quants can trade uncorrelated risk factors, isolating pure alpha.

\subsection{B) Model Problem Setup: PCA for Quantitative Finance}
Let $\mathbf{r} \in \mathbb{R}^n$ be a random vector representing the daily returns of $n$ equities. Let $\Sigma \in \mathbb{R}^{n \times n}$ be the covariance matrix of these returns, such that $\Sigma = \mathbb{E}[(\mathbf{r} - \boldsymbol{\mu})(\mathbf{r} - \boldsymbol{\mu})^T]$. Because $\Sigma$ is a symmetric, positive semi-definite matrix, the Spectral Theorem guarantees that it can be orthogonally diagonalized.

We define a new basis $\mathcal{B} = \{\mathbf{v}_1, \dots, \mathbf{v}_n\}$, where $\mathbf{v}_i$ are the orthonormal eigenvectors of $\Sigma$. The change-of-coordinates matrix from the eigenbasis $\mathcal{B}$ to the standard basis $\mathcal{E}$ is:
\[ P_{\mathcal{E} \leftarrow \mathcal{B}} = \begin{bmatrix} \mathbf{v}_1 & \mathbf{v}_2 & \dots & \mathbf{v}_n \end{bmatrix} = P \]
Because the basis is orthonormal, $P$ is an orthogonal matrix, meaning $P^{-1} = P^T$. Therefore, the change of basis from standard asset returns to uncorrelated principal portfolio returns is given by:
\[ [\mathbf{r}]_{\mathcal{B}} = P_{\mathcal{B} \leftarrow \mathcal{E}} [\mathbf{r}]_{\mathcal{E}} = P^T \mathbf{r} \]
Let $\mathbf{y} = P^T \mathbf{r}$. The covariance of the new risk factors $\mathbf{y}$ is:
\[ \text{Cov}(\mathbf{y}) = \mathbb{E}[\mathbf{y}\mathbf{y}^T] = \mathbb{E}[P^T \mathbf{r} \mathbf{r}^T P] = P^T \Sigma P = \Lambda \]
where $\Lambda$ is a diagonal matrix of eigenvalues. By applying a simple change of basis, the quantitative model transforms a highly correlated, complex state space into a set of completely uncorrelated, independent stochastic processes, allowing for independent risk modeling and optimal capital allocation.

% ==========================================
% PART 7: COMMON VARIATIONS
% ==========================================
\section{Part 7: Common Variations and Untested Concepts}

\subsection{Change of Basis for Linear Transformations}
The homework focuses on changing the basis for \textit{vectors}. A critical extension is changing the basis for \textit{matrices} (linear transformations). If $T: V \to V$ is a linear transformation, its matrix representation depends on the basis.
If $A$ is the matrix for $T$ in basis $\mathcal{B}$, and $C$ is the matrix for $T$ in basis $\mathcal{C}$, they are related by:
\[ C = P_{\mathcal{C} \leftarrow \mathcal{B}} A P_{\mathcal{B} \leftarrow \mathcal{C}} = P_{\mathcal{C} \leftarrow \mathcal{B}} A (P_{\mathcal{C} \leftarrow \mathcal{B}})^{-1} \]
Matrices $A$ and $C$ are called \textbf{similar matrices}.

\textbf{Example:} Let $T(\mathbf{x}) = \begin{bmatrix} 2 & 1 \\ 0 & 3 \end{bmatrix} \mathbf{x}$ in the standard basis $\mathcal{E}$. Find the matrix of $T$ relative to the basis $\mathcal{B} = \left\{ \begin{bmatrix} 1 \\ 0 \end{bmatrix}, \begin{bmatrix} 1 \\ 1 \end{bmatrix} \right\}$.
\textbf{Solution:} 
$P_{\mathcal{E} \leftarrow \mathcal{B}} = \begin{bmatrix} 1 & 1 \\ 0 & 1 \end{bmatrix}$. Its inverse is $P_{\mathcal{B} \leftarrow \mathcal{E}} = \begin{bmatrix} 1 & -1 \\ 0 & 1 \end{bmatrix}$.
The matrix in basis $\mathcal{B}$ is:
\[ [T]_{\mathcal{B}} = P_{\mathcal{B} \leftarrow \mathcal{E}} [T]_{\mathcal{E}} P_{\mathcal{E} \leftarrow \mathcal{B}} = \begin{bmatrix} 1 & -1 \\ 0 & 1 \end{bmatrix} \begin{bmatrix} 2 & 1 \\ 0 & 3 \end{bmatrix} \begin{bmatrix} 1 & 1 \\ 0 & 1 \end{bmatrix} = \begin{bmatrix} 2 & -2 \\ 0 & 3 \end{bmatrix} \begin{bmatrix} 1 & 1 \\ 0 & 1 \end{bmatrix} = \begin{bmatrix} 2 & 0 \\ 0 & 3 \end{bmatrix} \]
Notice the result is diagonal! This is because $\mathcal{B}$ is an eigenbasis for $T$.

% ==========================================
% PART 8: ADVANCED DIAGNOSTIC TESTING
% ==========================================
\section{Part 8: Advanced Diagnostic Testing: ``Find the Flaw''}

\textbf{Problem 1:} Let $\mathcal{B} = \left\{ \begin{bmatrix} 1 \\ 2 \end{bmatrix}, \begin{bmatrix} 3 \\ 4 \end{bmatrix} \right\}$ and $\mathcal{C} = \left\{ \begin{bmatrix} 0 \\ 1 \end{bmatrix}, \begin{bmatrix} 1 \\ 0 \end{bmatrix} \right\}$. Find $P_{\mathcal{C} \leftarrow \mathcal{B}}$.
\textit{Flawed Solution:} We set up the augmented matrix $[ \mathcal{B} \mid \mathcal{C} ]$:
\[ \begin{bmatrix} 1 & 3 & 0 & 1 \\ 2 & 4 & 1 & 0 \end{bmatrix} \sim \begin{bmatrix} 1 & 0 & 1.5 & -2 \\ 0 & 1 & -0.5 & 1 \end{bmatrix} \]
So $P_{\mathcal{C} \leftarrow \mathcal{B}} = \begin{bmatrix} 1.5 & -2 \\ -0.5 & 1 \end{bmatrix}$.
\textbf{Find the Flaw:} The student set up the augmented matrix backwards as $[ \mathcal{B} \mid \mathcal{C} ]$ instead of $[ \mathcal{C} \mid \mathcal{B} ]$.
\textbf{Correct Solution:} Set up $[ \mathcal{C} \mid \mathcal{B} ] = \begin{bmatrix} 0 & 1 & 1 & 3 \\ 1 & 0 & 2 & 4 \end{bmatrix}$. Swapping rows gives $\begin{bmatrix} 1 & 0 & 2 & 4 \\ 0 & 1 & 1 & 3 \end{bmatrix}$. Thus, $P_{\mathcal{C} \leftarrow \mathcal{B}} = \begin{bmatrix} 2 & 4 \\ 1 & 3 \end{bmatrix}$.

\vspace{0.5cm}

\textbf{Problem 2:} Given $P_{\mathcal{C} \leftarrow \mathcal{B}} = \begin{bmatrix} 2 & 1 \\ 5 & 3 \end{bmatrix}$ and $[\mathbf{x}]_{\mathcal{C}} = \begin{bmatrix} 1 \\ -1 \end{bmatrix}$, find $[\mathbf{x}]_{\mathcal{B}}$.
\textit{Flawed Solution:} We use the formula $[\mathbf{x}]_{\mathcal{B}} = P_{\mathcal{C} \leftarrow \mathcal{B}} [\mathbf{x}]_{\mathcal{C}}$.
\[ [\mathbf{x}]_{\mathcal{B}} = \begin{bmatrix} 2 & 1 \\ 5 & 3 \end{bmatrix} \begin{bmatrix} 1 \\ -1 \end{bmatrix} = \begin{bmatrix} 1 \\ 2 \end{bmatrix} \]
\textbf{Find the Flaw:} The student multiplied the $\mathcal{C}$-coordinate vector by $P_{\mathcal{C} \leftarrow \mathcal{B}}$, which requires a $\mathcal{B}$-coordinate vector as input.
\textbf{Correct Solution:} We must use the inverse matrix: $[\mathbf{x}]_{\mathcal{B}} = (P_{\mathcal{C} \leftarrow \mathcal{B}})^{-1} [\mathbf{x}]_{\mathcal{C}}$.
$(P_{\mathcal{C} \leftarrow \mathcal{B}})^{-1} = \frac{1}{6-5} \begin{bmatrix} 3 & -1 \\ -5 & 2 \end{bmatrix} = \begin{bmatrix} 3 & -1 \\ -5 & 2 \end{bmatrix}$.
$[\mathbf{x}]_{\mathcal{B}} = \begin{bmatrix} 3 & -1 \\ -5 & 2 \end{bmatrix} \begin{bmatrix} 1 \\ -1 \end{bmatrix} = \begin{bmatrix} 4 \\ -7 \end{bmatrix}$.

\vspace{0.5cm}

\textbf{Problem 3:} In $\mathbb{P}_2$, let $\mathcal{B} = \{1, t, t^2\}$ and $\mathcal{C} = \{t^2, t, 1\}$. Find $P_{\mathcal{C} \leftarrow \mathcal{B}}$.
\textit{Flawed Solution:} Since both are standard bases just written in different orders, the change of basis matrix is the identity matrix $I_3$.
\textbf{Find the Flaw:} The order of vectors in a basis strictly dictates the coordinate mapping, so changing the order changes the coordinate vectors.
\textbf{Correct Solution:} $[\mathbf{b}_1]_{\mathcal{C}} = [1]_{\mathcal{C}} = \begin{bmatrix} 0 \\ 0 \\ 1 \end{bmatrix}$. $[\mathbf{b}_2]_{\mathcal{C}} = [t]_{\mathcal{C}} = \begin{bmatrix} 0 \\ 1 \\ 0 \end{bmatrix}$. $[\mathbf{b}_3]_{\mathcal{C}} = [t^2]_{\mathcal{C}} = \begin{bmatrix} 1 \\ 0 \\ 0 \end{bmatrix}$.
$P_{\mathcal{C} \leftarrow \mathcal{B}} = \begin{bmatrix} 0 & 0 & 1 \\ 0 & 1 & 0 \\ 1 & 0 & 0 \end{bmatrix}$.

\vspace{0.5cm}

\textbf{Problem 4:} If $P_{\mathcal{C} \leftarrow \mathcal{B}} = \begin{bmatrix} 4 & 7 \\ 1 & 2 \end{bmatrix}$, find $P_{\mathcal{B} \leftarrow \mathcal{C}}$.
\textit{Flawed Solution:} To reverse the direction of the change of basis, we simply negate the matrix, treating it like a vector translation. $P_{\mathcal{B} \leftarrow \mathcal{C}} = -P_{\mathcal{C} \leftarrow \mathcal{B}} = \begin{bmatrix} -4 & -7 \\ -1 & -2 \end{bmatrix}$.
\textbf{Find the Flaw:} Change of basis is a linear transformation (multiplication), not a translation (addition), so the reverse operation is the matrix inverse, not the negative.
\textbf{Correct Solution:} $P_{\mathcal{B} \leftarrow \mathcal{C}} = (P_{\mathcal{C} \leftarrow \mathcal{B}})^{-1} = \frac{1}{8-7} \begin{bmatrix} 2 & -7 \\ -1 & 4 \end{bmatrix} = \begin{bmatrix} 2 & -7 \\ -1 & 4 \end{bmatrix}$.

\vspace{0.5cm}

\textbf{Problem 5:} Let $\mathcal{B} = \left\{ \begin{bmatrix} 1 \\ 1 \end{bmatrix}, \begin{bmatrix} 1 \\ 0 \end{bmatrix} \right\}$. Find the $\mathcal{B}$-coordinates of $\mathbf{x} = \begin{bmatrix} 3 \\ 2 \end{bmatrix}$.
\textit{Flawed Solution:} We use the dot product projection formula.
$c_1 = \frac{\mathbf{x} \cdot \mathbf{b}_1}{\mathbf{b}_1 \cdot \mathbf{b}_1} = \frac{5}{2}$. $c_2 = \frac{\mathbf{x} \cdot \mathbf{b}_2}{\mathbf{b}_2 \cdot \mathbf{b}_2} = \frac{3}{1} = 3$.
$[\mathbf{x}]_{\mathcal{B}} = \begin{bmatrix} 2.5 \\ 3 \end{bmatrix}$.
\textbf{Find the Flaw:} The dot product projection formula $c_i = \frac{\mathbf{x} \cdot \mathbf{b}_i}{\mathbf{b}_i \cdot \mathbf{b}_i}$ only works if the basis $\mathcal{B}$ is an orthogonal basis, which this is not ($\mathbf{b}_1 \cdot \mathbf{b}_2 = 1 \neq 0$).
\textbf{Correct Solution:} Solve the system $c_1\mathbf{b}_1 + c_2\mathbf{b}_2 = \mathbf{x}$.
$\begin{bmatrix} 1 & 1 \\ 1 & 0 \end{bmatrix} \begin{bmatrix} c_1 \\ c_2 \end{bmatrix} = \begin{bmatrix} 3 \\ 2 \end{bmatrix}$.
From the second row, $c_1 = 2$. From the first row, $2 + c_2 = 3 \implies c_2 = 1$.
$[\mathbf{x}]_{\mathcal{B}} = \begin{bmatrix} 2 \\ 1 \end{bmatrix}$.

\end{document}